%% how-to-use --- how the book is meant to be read, and what its boxes mean.
%%
%% Twin: frontmatter/en/how-to-use.tex. parity.py compares them token by
%% token, like every chapter.

\chapter*{\lblHowToUse}
\phantomsection
\addcontentsline{toc}{chapter}{\lblHowToUse}
\markboth{\lblHowToUse}{\lblHowToUse}

\noindent
Ta książka zakłada, że umiesz dodawać, mnożyć, policzyć procent i odczytać
wykres. Nie zakłada, że pamiętasz jakąkolwiek algebrę, i nigdy nie używa
symbolu, którego wcześniej nie wprowadziła. Każdy wzór jest też powiedziany
słowami, a tam, gdzie wzór tylko by przeszkadzał, wzoru nie ma. Jeśli
skończyłeś z matematyką w szkole i zawsze chciałeś wiedzieć, co ludzie mają
na myśli, kiedy mówią \enquote{chaos}, to jesteś czytelnikiem, dla którego ją
napisano.

Prosi cię za to o jedno: o udział. Książka jest zbudowana z małych kroków,
a na końcu wielu z nich pada pytanie.

\section*{Pętla}

Każdy rozdział czyta się tak samo:

\begin{enumerate}[leftmargin=1.6em, itemsep=3pt]
  \item \textbf{Przeczytaj krótki fragment.} Każdy podrozdział to kilka
        akapitów.
  \item \textbf{Zatrzymaj się i pomyśl.} Fioletowa ramka zadaje ci pytanie.
        Odpowiedz na nie \dash{} na kartce, na głos albo na zasadzie
        zgadywania \dash{} zanim pójdziesz dalej. Odpowiedź jest tuż pod
        spodem, więc zasłoń ją dłonią.
  \item \textbf{Uruchom kod}, jeśli masz obok komputer. Każdy listing to
        mały plik, a uruchomienie go to jedno polecenie.
  \item \textbf{Zrób laboratorium}, jeśli chcesz, żeby pomysł stał się
        twój: laboratorium to napisany do połowy program i test, który mówi
        ci, kiedy działa.
  \item \textbf{Rozwiąż sprawdzian} na końcu rozdziału i porównaj odpowiedzi
        z dodatkiem~\ref{app:answers}.
\end{enumerate}

\section*{Dlaczego książka wciąż zadaje ci pytania}

Pytanie zadane, zanim poznasz odpowiedź, jest niewygodne. Tak ma być.
Odpowiedź, po którą sięgnąłeś sam, nawet błędna, zostaje w głowie dużo
dłużej niż taka, którą tylko przeczytałeś, a błędny strzał, który potem
zostaje poprawiony, to najszybszy sposób, żeby odkryć, które twoje intuicje
cię zawiodą. W tym temacie jest ich pełno. Pierwszy strzał większości ludzi
co do tego, jak szybko rośnie mały błąd, albo co potrafi prosta reguła, jest
błędny, a książka jest ułożona tak, żebyś oddał ten strzał, zanim zobaczysz
dlaczego.

Więc kiedy ramka prosi, żebyś odpowiedział, zanim pójdziesz dalej, zrób to,
proszę. Nikt tego nie ocenia.

\section*{Laboratorium}

Wszystko w tej książce, co może sprawdzić komputer, zostało przez komputer
sprawdzone, a ty możesz to sprawdzić ponownie. Programy są w repozytorium
książki, w katalogu \texttt{code/}. Są napisane w Pythonie i potrzebują
jednego narzędzia, \pkg{uv}, które samo instaluje właściwą wersję Pythona i
wszystkie biblioteki. Kiedy już je masz, z katalogu \texttt{code/}
repozytorium:

\begin{shellcmd}
uv sync
uv run python ch01/falling_stone.py
uv run pytest -k l01_01
\end{shellcmd}

Pierwsza linia przygotowuje laboratorium, raz. Druga uruchamia listing z
rozdziału~\ref{ch:models}. Trzecia uruchamia test pierwszego laboratorium z
tego rozdziału. Nie musisz znać Pythona, żeby zacząć: pierwsze rozdziały
tłumaczą każdą linię, którą drukują, a programy przez całą książkę pozostają
krótkie.

\section*{Ramki}

Poza pytaniami tekst przerywają ramki, każda z jednym zadaniem:

\begin{description}[leftmargin=1.2em, itemsep=2pt]
  \item[\lblTrap.] Błąd, który popełnia prawie każdy, pokazany po tym, jak
    miałeś okazję popełnić go sam.
  \item[\lblWorld.] Prawdziwe miejsce, w którym pomysł się pojawia:
    dziedzina, urządzenie, praktyka.
  \item[\lblRigour.] Co książka stwierdza bez dowodu i gdzie ten dowód
    można znaleźć.
  \item[\lblNote, \lblWarning.] Dygresja i przestroga.
  \item[\lblProject.] Kawałek małej biblioteki chaosu tej książki, z której
    korzystają listingi.
  \item[\lblLab.] Coś do zrobienia przy klawiaturze, z testem, który mówi
    ci, kiedy to zrobiłeś.
\end{description}

\section*{Liczby, którym można ufać}

Każdą liczbę w tej książce, której nie da się policzyć w pamięci,
wyprodukował program, a książka drukuje to, co wypisał program, a nie liczbę,
którą ktoś przepisał. Proces budowania sprawdza każdą z nich ponownie przy
każdej zmianie. Tutaj ma to większe znaczenie niż zwykle, z powodu, któremu
książka poświęca cały rozdział: w układzie chaotycznym obliczenie może pomylić
się na ostatniej cyfrze i kilkadziesiąt kroków później być zupełnie błędne.
Rozdział~\ref{ch:computers} wyjaśnia, jak książka się przed tym chroni i
dlaczego niektóre jej liczby są celowo zaokrąglone.

\section*{Czytanie bez komputera}

Da się. Wynik każdego listingu jest wydrukowany obok niego, każdy wykres
jest na stronie, a pytania działają na papierze. Ominą cię tylko
laboratoria, a one są dla czytelników, którzy chcą pójść dalej, a nie po to,
żeby zamykać drogę do wywodu.

\section*{Dwa języki, jedna książka}

Książka ukazuje się po angielsku i po polsku z jednego źródła. Oba wydania
zawierają te same rozdziały, podrozdziały, wykresy, liczby i pytania w tej
samej kolejności, a proces budowania porównuje je, żeby tego dopilnować.
